\chapter*{Abstract}
\addcontentsline{toc}{chapter}{Abstract}

Distributed Denial-of-Service (DDoS) attacks against multi-tenant data-center
networks continue to grow in volume and subtlety, yet the dominant defensive
posture remains \emph{reactive}: mitigation fires only after the attack
has crossed a static threshold. Reactive mitigation inherits a hard lower
bound on response time equal to the virtual-network-function (VNF) boot
latency --- in our measurements, approximately six seconds --- during which
the attack is already draining tenant bandwidth.

This thesis proposes \padonap{}, an AI-augmented NFV orchestration pipeline
that integrates (i) a 17-feature entropy/rate telemetry stage, (ii) an XGBoost
seven-class detector with SHAP interpretability, (iii) a
Transformer+LSTM four-horizon forecaster (\(t{+}30{\dots}120\) s), (iv) a
five-tier graduated policy engine with hysteresis, (v) a linear-programming
multi-tenant SLA allocator, and (vi) an ONAP MANO-compliant control loop
(SO, CLAMP, Policy Framework, DMaaP).

Evaluated on CICDDoS2019 with temporal 80/20 split and real-BENIGN
oversampling, the detector attains accuracy \(0.9825\), macro-F1 \(0.9642\)
and AUC \(0.9958\); the forecaster's horizon-AUC degrades
monotonically from \(0.988\) at \(t{+}30\)s to \(0.971\) at \(t{+}120\)s,
which is consistent with the physical forecast-horizon trade-off.
On an eight-scenario orchestration benchmark \padonap{} passes 8/8 scenarios,
reaches \tier{2} pre-positioning in \(505\) ms versus \(6006\) ms for
\tier{3} reactive mitigation (\(10.9\times\) speed-up), and achieves a
\(65\) s lead-time advantage over a rule-based threshold baseline on
gradual SYN-ramp attacks. Under out-of-distribution (OOD) traffic the
threshold baseline over-escalates and FAILs two of eight scenarios, whereas
\padonap{} degrades gracefully to \tier{0}--\tier{1}. URLLC-tenant SLA floors
(\(150\) Mbps) are preserved throughout.

The primary novelty is architectural: \padonap{} is, to our knowledge, the
first system combining ML-based multi-horizon forecasting, ONAP-compliant
lifecycle management, and LP-based multi-tenant SLA guarantees within a
single closed-loop orchestration pipeline.

\vspace{1em}
\noindent\textbf{Keywords:} DDoS mitigation; NFV orchestration; ONAP;
proactive defense; DDoS forecasting; XGBoost; Transformer--LSTM; SLA;
data-center networks.

\clearpage

\chapter*{Tóm tắt}
\addcontentsline{toc}{chapter}{Tóm tắt}
\selectlanguage{vietnamese}

Tấn công từ chối dịch vụ phân tán (DDoS) nhằm vào mạng trung tâm dữ liệu
đa thuê bao ngày càng lớn về quy mô và tinh vi về kỹ thuật, trong khi
các cơ chế phòng thủ chủ đạo vẫn mang tính \emph{phản ứng}: việc giảm thiểu
chỉ được kích hoạt sau khi lưu lượng đã vượt ngưỡng tĩnh, dẫn tới độ trễ
đáp ứng tối thiểu bằng thời gian khởi động VNF --- khoảng sáu giây trong
thực nghiệm --- giai đoạn mà tấn công đã làm suy giảm băng thông thuê bao.

Luận văn đề xuất \padonap{}, một pipeline điều phối NFV tăng cường AI tích
hợp: (i) thu thập 17 đặc trưng entropy/rate, (ii) bộ phát hiện XGBoost bảy
lớp giải thích được bằng SHAP, (iii) bộ dự báo Transformer+LSTM bốn tầm
dự báo (\(t{+}30{\dots}120\) s), (iv) bộ máy chính sách năm mức có
hysteresis, (v) bộ phân bổ băng thông SLA đa thuê bao bằng quy hoạch tuyến
tính, và (vi) vòng điều khiển tuân thủ ONAP MANO.

Trên CICDDoS2019 với chia tập thời gian 80/20, bộ phát hiện đạt độ chính
xác \(0{,}9825\), macro-F1 \(0{,}9642\), AUC \(0{,}9958\); AUC dự báo suy
giảm đơn điệu từ \(0{,}988\) (\(t{+}30\)s) xuống \(0{,}971\) (\(t{+}120\)s).
Trên 8 kịch bản đánh giá, \padonap{} PASS toàn bộ 8/8, đạt mức
\tier{2} chủ động trong \(505\) ms so với \(6006\) ms của \tier{3} phản ứng
(nhanh hơn \(10{,}9\) lần), và tạo độ lợi thời gian \(65\) giây so với
đường ngưỡng không-AI trên tấn công SYN tăng dần. Với lưu lượng
ngoài-phân-phối (OOD), đường ngưỡng FAIL 2/8 kịch bản trong khi \padonap{}
suy giảm có kiểm soát. Sàn SLA URLLC \(150\) Mbps được bảo toàn.

\noindent\textbf{Từ khóa:} giảm thiểu DDoS; điều phối NFV; ONAP; phòng thủ
chủ động; dự báo DDoS; XGBoost; Transformer--LSTM; SLA; mạng trung tâm
dữ liệu.

\selectlanguage{english}
